\documentclass[10pt,a4paper]{amsart}
\usepackage[margin=19mm]{geometry}
\usepackage{amsmath,amssymb,mathtools}
\usepackage[scr=rsfs,cal=euler]{mathalfa}
\usepackage{xfrac}
\usepackage[unicode,hidelinks,bookmarks=false]{hyperref}
\newcommand{\ie}{i.e.\ }
\newcommand{\cf}{cf.\ }
\newcommand{\resp}{respectively\ }
\newcommand{\Subsection}{\subsection}
\providecommand{\nobreakdash}{\nobreak\hbox{-}}
\setlength{\emergencystretch}{2em}
\multlinegap0pt
\newtheorem{lemma}{Lemma}
\title{A note on an inequality involving the sides and medians in a triangle}
\author{Peter Vassilev}
\date{Source: arXiv:2602.15054v2 (CC BY 4.0)}
\begin{document}
\maketitle

\noindent\emph{Source and licence.} A note on an inequality involving the sides and medians in a triangle. Version arXiv:2602.15054v2 (CC BY 4.0), distributed by arXiv under the Creative Commons Attribution 4.0 International licence, http://creativecommons.org/licenses/by/4.0/, as recorded in the arXiv licence metadata for this exact version and shown on the abstract page https://arxiv.org/abs/2602.15054v2. Full licence notice: \texttt{licenses/arxiv-2602.15054-CC-BY-4.0.txt}.\par\smallskip

\noindent\textit{Editorial scope.} Counted unit: the article's lemma isosceles (source0.tex lines 59--61). The article's complete original proof of it is delivered with it (source0.tex lines 62--109, one \texttt{proof} environment of 48 lines). No mathematics is written, repaired or re-derived: the statement and the proof are the article's own lines, in the article's own notation, and no retained line is substituted or re-flowed.  Retained uncounted as the source-local context the proof needs: lines 34--36.  Nothing was omitted from any retained span, so the package carries no omission record.  No retained line contains a citation, so the wrapper carries no bibliography and no bibliography entry.  No retained line contains a figure environment, an image-inclusion command or drawing code, so no image is delivered and the package declares no asset dependency.  Editorial formatting only, in addition: an AMS \texttt{amsart} preamble that restates the article's own declarations and macros used by the retained lines, namely the environments \texttt{lemma} on separate counters, together with the standard packages those lines need.  The retained environments print in this document as Lemma 1 in order of appearance, whereas the article prints the same items with the numbering of its own full document.  The complete native source of the article as distributed in the arXiv e-print for this exact version is bundled unchanged as \texttt{sources/194768/source0.tex}.
\begin{equation}\label{INEQ}
 \sqrt{bc} m_a+ \sqrt{ac} m_b+\sqrt{ab} m_c \geq a m_a+ b m_b+c m_c 
\end{equation}

\begin{lemma}\label{isosceles} 
Inequality \eqref{INEQ} is true for any isosceles triangle.
\end{lemma}

\begin{proof} First we observe that \eqref{INEQ} turns into identity for an equilateral triangle. 

Thus, without loss of generality we can have either: 1)$a=b<c$ or 2)$a<b=c.$

Let us consider case 1). Then \eqref{INEQ} becomes
\begin{equation}\label{INEQFIRSTISOS}
2\sqrt{ac} m_a+a m_c \geq 2 a m_a+c m_c 
\end{equation}
Dividing \eqref{INEQFIRSTISOS} by $c^2,$ and putting $x=\frac{a}{c}<1,$ it is equivalent to:
\begin{align*}
2 \sqrt{x} \sqrt{2+x^2}+x\sqrt{4x^2-1} - 2x \sqrt{2+x^2}- \sqrt{4x^2-1} \\
= 2\sqrt{x}(1-\sqrt{x})\sqrt{2+x^2}+(\sqrt{x}-1)(\sqrt{x}+1)\sqrt{4x^2-1} \\
=(1-\sqrt{x})(2\sqrt{2x+x^3}-(\sqrt{x}+1)\sqrt{4x^2-1}) \geq 0
\end{align*}
But we have,
\begin{align*}
2\sqrt{2x+x^3}-(\sqrt{x}+1)\sqrt{4x^2-1} \geq 2(\sqrt{2x+x^3}-\sqrt{4x^2-1})\\
=\frac{2({2x+x^3}-{4x^2+1})}{\sqrt{2x+x^3}+\sqrt{4x^2-1}}=\frac{2(1-x)(-x^2+3x+1)}{\sqrt{2x+x^3}+\sqrt{4x^2-1}} \geq 0
\end{align*}
since $1>x.$ Hence, the desired inequality is fulfilled.



Let us consider case 2). Then \eqref{INEQ} becomes
\begin{equation}\label{INEQSECONDISOS}
c  m_a+2\sqrt{ac} m_c \geq a m_a+2 c m_c 
\end{equation}
Putting $x=\frac{a}{c}<1,$ after division by $c^2,$ \eqref{INEQSECONDISOS} is equivalent to:
\[
\sqrt{4-x^2}+2\sqrt{x}\sqrt{1+2x^2} \geq x\sqrt{4-x^2}+2\sqrt{1+2x^2}
\]
Hence, 
\[
(1- x)\sqrt{4-x^2}-2\sqrt{1+2x^2}(1-\sqrt{x})\geq 0
\]
Simplifying,
\[
(1-\sqrt{x})\left((1+\sqrt{x})\sqrt{4-x^2}-2\sqrt{1+2x^2}\right)\geq 0
\]
Using the fact that $1+\sqrt{x} \geq 1+x,$ we have:
\begin{align*}
(1-\sqrt{x})\left((1+\sqrt{x})\sqrt{4-x^2}-2\sqrt{1+2x^2}\right)\\
\geq (1-\sqrt{x})\left((1+x)\sqrt{4-x^2}-2\sqrt{1+2x^2}\right)\\
=\frac{(1-\sqrt{x})\left((1+x)^2(4-x^2)-4(1+2x^2)\right)}{(1+x)\sqrt{4-x^2}+2\sqrt{1+2x^2}} \\
=\frac{(1-\sqrt{x})x\left(8-5x-2x^2-x^3\right)}{(1+x)\sqrt{4-x^2}+2\sqrt{1+2x^2}}  \geq 0,
\end{align*}
since $x<1.$
\end{proof}

\end{document}
